\documentclass[a4paper,12pt]{article}
\usepackage{amsmath,amssymb,amsthm}
\usepackage[utf8]{inputenc}

\title{拡張連続性公理（ECA）における Solver集合の定義と公理選択空間 $I$ の形式定義}
\author{}
\date{}

\begin{document}
\maketitle

\section*{定義 1（Solver 集合）}
\[
S \quad(\text{Solver集合})
\]
を  \emph{全公理集合}（選択可能な全ての公理の全集合）とする。  
例：巨大基数公理、ZFCなどの既にある公理系など

これは理論内で選択し得る公理や公理スキームをすべての要素として含む集合のことである。

\section*{定義 2（公理選択空間 $I$）}
公理選択空間 $I$ を次で定義する：
\[
I \subseteq \mathcal{P}(S).
\]
すなわち、$I$ の各元 $i\in I$ は $S$ の部分集合であり「ある公理選択の結果として採用される公理集合」を表す。

各 $i\in I$ は有限または可算な要素数 $N(i)\in\mathbb{N}\cup\{\aleph_0\}$ を持ち、特に要素数が $0$ の場合は空集合とする。

\section*{定義 3（分類：離散・連続・中間）}
公理選択空間 $I$ の元 $i\in I$ に対して，次の三つのクラスのいずれかに属するものとする（以下は記法上の分類である）：
\begin{itemize}
  \item \(\mathcal{C}_{\mathrm{disc}}\)（離散的公理集合の族）: 個々の要素が明確に分離され得る公理集合群。
  \item \(\mathcal{C}_{\mathrm{cont}}\)（連続的公理集合の族）: 連続性や連続的変形が意味を持つ範疇として扱う公理集合群。
  \item \(\mathcal{C}_{\mathrm{mid}}\)（中間的公理集合の族）: 前二者の中間に位置するスペクトラム的な集合群。
\end{itemize}
これらは互いに排反であるとは限らないが、任意の $i\in I$ に対して少なくとも一つまたは複数の集合が割り当てられるものとする（分類の技術的条件は別途定式化する）。

\section*{定義 4（公理数と空集合の扱い）}
任意の $i\in I$ に対し、その有する公理の個数を $N(i)$ と表す。  
特に $N(i)=0$ の場合は$i=\varnothing$（空集合）である。

命題論理ないし数学命題 $\varphi$ に対して，$\varphi$ の解決（あるいは記述）に必要とされる公理集合を
\[
A_\varphi \subseteq S
\]
と定める。
もし $A_\varphi \neq \varnothing$ であれば $A_\varphi \in I$（または $I$ に包含される代表元）として取り扱う。  
一方 $A_\varphi = \varnothing$（すなわち「その命題に対して指定された公理が存在しない／選ばれていない」）である場合には、まず空集合を \emph{代表元} として追加し、その代表元に属する公理群の要素をいずれかの集合 $\mathcal{C}_{\mathrm{disc}},\mathcal{C}_{\mathrm{cont}},\mathcal{C}_{\mathrm{mid}}$ またはその中の複数からなる集合から要素を選び、それを用いて $\varphi$ を求める（この操作の詳細アルゴリズムは ECA の別節にて定式化する）。

\section*{補足}
\begin{itemize}
  \item Solver集合の定義は集合論や測度論を用いる
  \item 分類（離散／連続／中間）の厳密基準（例えば測度論的基準，位相的基準など）は別途定義する。
  \item $I$ は $S$ から導出される（すなわち $I\subseteq\mathcal{P}(S)$）が、$I$ は「公理選択空間」という写像に置き換えるため、厳密にはSで得た集合を測度に変換しそれを写像として扱うこととなる。
\end{itemize}

\section{ここまでの内容}
定義：Solver集合（数学的基礎）

\subsection{定義1（Solver集合）}
Sを全公理集合とし、次を満たす集合とする：

1.S は集合とし $S \in \mathbf{Set}$ となるものとする。
2. 集合の各要素 $s \in S$ は「公理（または公理スキーム）」を表す ($s : \text{axiom or axiom scheme}$) 。
3. S は任意の数学体系で使用し得る公理をすべて含む (例1を参照)。

例1：
\begin{itemize}
  \item ZFC
  \item 大きな基数公理 (巨大基数)
  \item 各種集合論的独立命題
  \item 正則性・選択・構成可能性に関する公理群
  \item 体系間の差異を生む公理スキーム
\end{itemize}

等々。

4. S は可算とは限らず、一般には巨大基数レベルの大きさを持ちうる。
すなわち $|S| \quad \text{は任意の基数になり得る}$


---

\subsection{定義 2（測度論的前提）}

Solver集合Sは後に測度へ置き換えられるように、以下の性質を満たす必要がある：

(1) S の冪集合に対して σ-代数を定義可能であること。

すなわち $\mathcal{F} \subseteq \mathcal{P}(S)$ であり、これは「空集合と S を含むもの」かつ「補集合が閉集合であること」、そして「可算和に関しても閉じた集合族」を取れるだけの一般性があることを示す。

\subsubsection{補足}
この段階では σ-代数そのものを明示しなくても良い


(2) S 上に測度 $\mu$ を定義可能な構造であること。

この測度を表す式は $\mu : \mathcal{F} \to [0,1]$ である。
ただし、ここで用いる $\mu$ はECA の I 空間で使われる $\mu$ と同名だが「数学的な概念としては別物」であることに注意すること。
Solver集合の段階では、単に測度論を用いる際の準備を整えておくために記述されるものである。



\subsection{σ-代数（sigma-algebra）の定義}

集合 $X$ 上の族 $\mathcal{F} \subseteq \mathcal{P}(X)$ が σ-代数であるとは、次を満たすときのことを指す。

1. 空集合と全体集合を含む
$\varnothing \in \mathcal{F}, \quad X \in \mathcal{F}$

2. 補集合が閉じている (閉集合)
$A \in \mathcal{F} \Rightarrow X\setminus A \in \mathcal{F}$

3. 可算和（可算個の unión）が閉じている (閉集合)
$A_1, A_2, A_3, \dots \in \mathcal{F} \quad \Rightarrow \quad \bigcup_{n=1}^{\infty} A_n \in \mathcal{F}$

このうちの3を “完全（可算）加法性” と呼ぶ。

σ-代数は可算個の和（union）に対して閉じており、また可算個の積（intersection）にも閉じている (ド・モルガンの法則で示せる) 。
つまり有限操作だけではなく「可算操作」に対しても閉じているという意味で完全に（可算）加法的な族であることからそう呼ばれる。

\subsection{Solver集合から測度へ}
Solver集合 $S$ を測度として置き換えるなら、σ-代数 $\mathcal{F} \subseteq \mathcal{P}(S)$ を取ることで ($S, \mathcal{F},  \mu$) という測度空間を定義可能になる。

\subsection{より深く議論できる点}
・Solver集合 $S$ の σ-代数は何を取るべきか？
・冪集合 $\mathcal{P}(S)$ をそのまま σ-代数にするべきか？
・S が巨大基数レベルなら σ-代数の性質はどうなる？


\subsection{}
なぜ Solver集合に σ-代数を導入するのが合理的なのか？

Solver集合 $S$ は「全公理集合の中から選び取られた公理」を要素として持つ集合である。
つまり $S$ の各元は「公理選択をしている状態」と捉えることができる。
その集合の部分集合は「公理選択によって選択された公理の集合（性質・条件）」と考えることができる。

これらを測度論で扱う場合、当然 $\mathcal{F} \subseteq \mathcal{P}(S)$ という σ-代数を導入するのが最も一般的かつ抽象的に正しい。

その理由は次の通りである。

1. 公理選択の状態集合に対して“可算操作で閉じた体系”が必要
Solver集合では「条件集合」や「分類」などを扱うため、それらが可算操作で閉じている必要がある。
したがって、σ-代数がそれに対応する。

2. 測度 μ を導入できる唯一の自然な道具
測度論では $\mu: \mathcal{F} \to [0,1]$ を定義するには、定義域が σ-代数でなければいけない。

この理論では μ(x) を [0,1] の連続値（小数点 3 桁、微分可能）で扱うため、次のような利点も生まれる。
・μ を測度として扱う
・μ を導関数（勾配）まで考えられる
・μ を Solver集合上の連続構造に関係させられる


3. 公理選択空間 I を「確率空間」に近い形で見られる
I の本質は「公理選択された後の空間」であり、Solver集合 $S$ は「公理選択を行う前の全体を表す集合」と考えることができ、すると自然に $(S, \mathcal{F}, \mu)$という測度空間が生まれ、公理選択空間I は μ の上での「ある種の事象（イベント）」として扱える。
したがって言えば、I の性質を解析する数学的土台として σ-代数は非常に相性が良い。


また、次のようなことも考慮する価値があるといえる。

・Solver集合の σ-代数を 冪集合全体にするべきか
・生成される σ-代数（例えば有限生成）にするべきか
・I に自然に誘導される測度 μ の形
・μ の微分可能性をどう定義するか（Radon–Nikodym を使うか）


\end{document}
